\section{Broader Impacts}
\label{sec:broader}

\subsection{Education and Mentoring}

The PIs have an outstanding record in education and curricular
development. PI Kruegel and PI Vigna teach several classes related to
the proposed research, including security, programming languages,
network applications, and software engineering. These courses are
routinely among the highest-ranked in their departments, and the PIs
have received awards for their teaching. For instance, PI Vigna
received the UCSB Academic Senate Distinguished Teaching Award in
2004---the highest teaching award at UC Santa Barbara, given annually
to the four best professors across the university. Both PI Kruegel
and PI Vigna have received the College of Engineering Student Council
Outstanding Professor Award, given annually to the ``Best Professor in
Computer Science.'' In 2025, PI Kruegel received the UCSB Outstanding
Graduate Mentor Award, given to three UCSB faculty members for
exceptional mentoring.

The PIs plan to introduce modules on \cps security explicitly into
their graduate and undergraduate classes, including hands-on lab
exercises in which students find and exploit bugs in \cps components
and propose remediations. The \sysname framework itself, which is
designed to run continuously against reproducible twin environments,
is particularly well-suited to classroom use: students can be given
their own twin instances in which to experiment without any risk to
real equipment.

The PIs have a history of including undergraduate and high-school
students in their research, with students often listed as co-authors
on publications submitted to international conferences. For example,
over the last few years, the PIs have published several papers with
Audrey Dutcher, Paul Grosen, and Jessie Grosen---all of whom joined
the lab as high-school interns, began their respective research
projects, and later continued contributing as undergraduates at UCSB.

Additionally, the PIs have leveraged prior NSF REU funding to support
several undergraduate students, including April Sanchez (a URM female
student) in 2021 and Gabriel Pizarro in 2024--2026. April was
involved in research on developing novel security competitions,
graduated in 2022 with distinction, and received the Computer Science
Student of the Year Award. Gabriel Pizarro worked on using planning
techniques for composing vulnerabilities and, in 2025, received the
Honorable Mention in the Computing Research Association Outstanding
Undergraduate Researcher Award. The PIs also run a summer internship
program that hosts approximately ten national and international
undergraduate students each year for three to six months; over the
past decade, the program has hosted more than 100 interns. The PIs
will continue to run this internship program and plan to involve
students in each of the project's thrusts.

The PIs will make a focused effort to include underserved students in
undergraduate research. To this end, PIs Kruegel and Vigna will draw
on their experience as faculty mentors in the Early Research Scholars
Program (ERSP), an NSF-supported, year-long research apprenticeship
program in the UCSB CS department designed to support students in
their first research experience. The PIs plan to work with a team of
three to four ERSP students on developing sophisticated firmware
emulators to analyze the security of monolithic firmware. As part of
this effort, students will study various emulation approaches and
design novel techniques to gather information about the hardware
configuration of \cps components. A letter of collaboration from
Ziad Matni, who leads the ERSP program, is attached to this proposal.

\subsection{Community Outreach}

The PIs will organize a series of security competitions centered on
\cps security. The PIs have participated with their hacking team
(Shellphish) in numerous capture-the-flag (CTF) competitions and
created the first distributed educational CTF competition, the iCTF.
The iCTF has run every year since 2003, and in 2026 introduced a new
type of challenge that requires developing agents to solve security
problems. From this vantage point, the PIs have seen firsthand the
effect that non-traditional educational tools and cybersecurity
competitions have on participants, who invest substantial resources
in preparing, executing, and post-processing their participation.

As part of this project, the PIs will organize an annual CTF
competition (separate from the iCTF) focused specifically on \cps
security. This competition will raise awareness, develop participant
capabilities, and motivate students to look more deeply into \CPSs
and their vulnerabilities. The competition will be open to the public,
with challenges at a range of difficulty levels, specifically to
attract students with diverse backgrounds and skill levels. The PIs
will leverage the work of ERSP students, whose focus is on firmware
emulation, to develop \cps components for these live competitions.

\subsection{Technology Transfer, Open Source, and Dissemination}

The PIs have a long history of making their research artifacts publicly
available. This includes Anubis and Wepawet, two systems hosted at
UCSB that analyzed tens of millions of malware programs and malicious
URLs, and \texttt{angr}, a binary analysis framework with tens of
thousands of downloads and (at the time of writing) more than 8{,}200
stars on GitHub. The PIs plan to release the \sysname framework as
open source under a permissive license. The goal is not only to
disseminate the specific tooling developed in this project but also to
lower the entry barrier for subsequent research on continuous
assessment of cyberphysical systems.

The PIs have a successful track record of working with industry
partners. PIs Kruegel and Vigna founded Lastline, a network security
and malware analysis company that grew to approximately 150 employees
before it was acquired by VMware in 2020. Industry collaborations
provide active avenues for support and the mutual exchange of ideas,
and allow the PIs to transfer research results and software prototypes
to practice. Interaction with industry partners through internships
and research meetings will strengthen the connection between our
basic research and the requirements of real-world deployments.

For scholarly dissemination, results will be published in top-tier
security venues, including the IEEE Symposium on Security and Privacy,
the ACM Conference on Computer and Communications Security (CCS), the
USENIX Security Symposium, and the Network and Distributed Systems
Security Symposium (NDSS).
